\documentclass[11pt, letterpaper]{article}

\input{/home/hyperion/.config/LatexHub/preambles/base.tex}
\input{/home/hyperion/.config/LatexHub/preambles/tikz.tex}
\input{/home/hyperion/.config/LatexHub/preambles/diagrams.tex}
\input{/home/hyperion/.config/LatexHub/preambles/macros-cat-theory.tex}
\input{/home/hyperion/.config/LatexHub/preambles/macros-algebra.tex}
\input{/home/hyperion/.config/LatexHub/preambles/macros-geometry.tex}
\input{/home/hyperion/.config/LatexHub/preambles/homework-formatting-global.tex}
\input{/home/hyperion/.config/LatexHub/preambles/refs-basic.tex}
\crefname{problem}{problem}{problems}
\Crefname{problem}{Problem}{Problems}
\usepackage{titlepageBU}
\title{Homework 1}
\courseID{MA745}
\begin{document}
\makeproblem

\begin{problem}\label{prob:1}
	Let $X$ be an affine variety. Show that the coordinate ring $A(X)$ is a field if and only if $X$ is a single point.
\end{problem}
\begin{proof}
	($\implies $) Let $X=\left\{ p \right\} \subseteq \mathbb{A}_k^n$ be a point, and write $p=(p_1, \ldots , p_n)$. Note that $p$ is precisely the vanishing locus of the ideal $J = (x_1 - p_1, \ldots , x_n - p_n)$. Note that this is clearly a maximal ideal, since $k[x_1, \ldots , x_n] / J \cong k$ is a field. By Hilbert's Nullstellensatz, $A(X) = k[x_1, \ldots , x_n] / I(V(J)) \cong k[x_1, \ldots , x_n] / \sqrt{J} $. Since all maximal ideals are radical, $\sqrt{J} =J$. Hence, $A(X) \cong k$.

	($\impliedby $) Now let $X$ be an affine variety such that $A(X)$ is a field. Then $I(X)$ is necessarily a maximal ideal. By the above argument, any point $p\in X$ yields an ideal $I(p) = (x_1 - p_1, \ldots , x_n - p_n)$. Since Hilbert's Nullstellensatz is inclusion-reversing and $\left\{ p \right\} \subseteq X$, we have $I(X) \subseteq I(p)$. Since $I(X)$ is maximal, we have $I(X) = I(p)$. Since $V(I(X)) = X$, we have $X = V(I(p)) = \left\{ p \right\} $ by the Nullstellensatz, hence $X$ is a point.
\end{proof}



\begin{problem}
	Determine the radical of the ideal $J=\left\langle x^3_1 - x^6_2, x_1x_2-x_2^3 \right\rangle \subseteq \mathbb{C} [x_1,x_2]$. (Hint: Hilbert's Nullstellensatz may be useful here.)
\end{problem}
\begin{proof}
	Sketch of the idea: by Hilbert's Nullstellensatz, $\sqrt{J} =I(V(J))$. The vanishing locus of this ideal should be computable. The radical is then just all functions that vanish on this locus. The first generator probably give something like $x_1 - x_2^2$, and idk what the right is.



	Note that $x_1x_2 - x_2^3 = x_2(x_1 - x_2^2)$. 
\end{proof}
\begin{problem}
	Let $X\subseteq \mathbb{A} _k^3$ be the union of the three coordinate axes. Compute generators for the ideal $I(X)$, and show that $I(X)$ cannot be generated by fewer than 3 elements.
\end{problem}
\begin{proof}
	A useful result might be Lemma 1.4 in Galthmann's notes, which says that $V(S_1)\cup V(S_2) = V(S_1S_2)$. This implies $X = V(xyz)$, and the Nullstellensatz says $I(X) = \sqrt{\left\langle xyz \right\rangle } $.


	Galthmann Lemma 1.4 says that $V(S_1) \cup V(S_2) = V(S_1S_2)$. Since the $x_i$-axis is precisely $V(x_i)$, we see that $X = V(xyz)$. Thus $f\in I(X)$ if and only if $xyx=0$ implies $f(x,y,z)=0$.

	Since $k$ is a field, $k[x,y,z]$ is an integral domain, hence it has no zero divisors. Thus $xyz=0$ implies at least one of $x,y,z$ is 0. We claim $I(X)=\left\langle x,y,z \right\rangle $. It suffices to show that all terms of some $f\in I(X)$ are of the form $\lambda x^ny^mz^{\ell } $, since such a term is necessarily in $\left\langle x,y,z \right\rangle $. Write
	\[
		f(x,y,z) = \sum_{n,m,\ell \geq 0} f_{n,m,\ell } x^ny^mz^{\ell } ,
	\]
	where $f_{n,m,\ell } \in k$ are the coefficients of $f$.

	Recall that ideals of a ring $R$ are precisely $R$-submodules of $R$. Hence, to show $I(X)$ cannot be generated by less than 3 elements, it suffices to show that $\left\{ x,y,z \right\} $ is a basis for the free module $I(X)$. Since we have shown that they span the space, it suffices to show $k[x,y,z]$-linear independence. Since there are no relations between $x,y,z$, this is so trivial that it's difficult to come up with a proof which rigorously expresses this fact. I came up with the following. Recall that
	\[
		k[x] \otimes k[y] \otimes k[z] \cong k[x,y,z]
	\]
	via the identification $x \otimes y \otimes z\mapsto xyz$. Hence we can equivalently show that the elements $x\otimes 1 \otimes 1$, $1\otimes y\otimes 1$, and $1\otimes 1\otimes z$ are linearly independent.
\end{proof}
\begin{problem}
	Let $Y\subseteq \mathbb{A} _k^n$ be an affine variety, and let $\pi : k[x_1, \ldots , x_n] \to A(Y)$ be the quotient map.
	\begin{enumerate}
		\item Show that $V_Y(J)=V(\pi ^{-1} (J))$ for every ideal $J\subseteq A(Y)$.
		\item Show that $\pi ^{-1} (I_Y(X)) = I(X)$ for every affine subvariety $X\subseteq Y$.
		\item Use (a) and (b) to deduce the interesting part of the relative Nullstellensatz $I_Y(V_Y(J)) \subseteq \sqrt{J} $ for every ideal $J\subseteq A(Y)$ from the corresponding absolute statement $I(V(J))\subseteq \sqrt{J} $ for every ideal $J\subseteq k[x_1, \ldots , x_n]$. In particular, conclude that there is an inclusion-reversing bijection between affine subvarieties of $Y$ and radical ideals of $A(Y)$.
	\end{enumerate}
\end{problem}

\end{document}
